\documentclass[a4paper,11pt,oneside]{book}

%========================================================================%
\usepackage{courier}
\usepackage[utf8]{inputenc}   
\usepackage[T1]{fontenc}      

\usepackage{graphicx}
\usepackage{xstring}
\usepackage{etoolbox}
\usepackage{float}
\usepackage{tabularx}
\usepackage[font=small,skip=2pt]{caption}
\usepackage{subcaption}
\captionsetup{within=none}
\usepackage{wrapfig}
\usepackage{multicol}
\usepackage{eurosym}
\usepackage{footmisc}
\usepackage{hyperref}
\usepackage{hyperxmp}
\usepackage{booktabs}
\usepackage[toc,page,title,titletoc]{appendix} % úsalo solo si agregas apéndices
\renewcommand*\appendixpagename{Apéndices}

\usepackage[spanish]{babel}
\addto\captionsspanish{\renewcommand{\tablename}{Tabla}}
\addto\captionsspanish{\renewcommand{\listtablename}{Índice de Tablas}}

% TikZ se carga DESPUÉS de babel: el español vuelve activos los símbolos < y >,
% y la librería 'babel' de TikZ es la que evita que eso rompa las flechas.
\usepackage{tikz}
\usetikzlibrary{arrows.meta,babel}

\usepackage[dvipsnames]{xcolor}
\definecolor{azul}{RGB}{81, 131, 187}   % Azul institucional UNIR
\def\azul{\color{azul}}
\def\an{\bfseries\azul}

\usepackage{geometry}
\geometry{
    paper=a4paper,
    left=3.5cm,
    right=1.5cm,
    top=2.5cm,
    bottom=2.5cm,
    headheight=0.75cm,
    headsep=1.5cm,
    footskip=1.5cm,
}

\usepackage{lastpage}
\usepackage{fancyhdr}
\pagestyle{fancy}
\patchcmd{\chapter}{\thispagestyle{plain}}{\thispagestyle{fancy}}{}{}

\usepackage{array}
\usepackage{setspace}
% La plantilla trae \setstretch{1.5}. El profesor fijó un máximo de 15
% hojas y con 1,5 el documento ocupa 17 sin que sobre contenido. Se baja a
% 1,25, que sigue siendo interlineado amplio y no altera la estructura
% obligatoria (introducción, desarrollo, resultados, discusión,
% conclusiones y bibliografía) ni ningún otro elemento de la plantilla.
% Para volver al valor original, cambiar 1.25 por 1.5.
\setstretch{1.25}
\usepackage[autostyle=true]{csquotes}

% Bibliografía en formato APA (7.ª ed.) real, vía el paquete biblatex-apa:
% sangría francesa, orden alfabético autor-año, "et al." según reglas APA.
\usepackage[backend=biber,style=apa,sorting=nyt]{biblatex}
\addbibresource{referencias.bib}
\renewcommand*{\cite}[1]{\parencite{#1}}
\newcommand*{\apa}[1]{\citeauthor{#1} (\citeyear{#1})}
\usepackage{ifxetex,ifluatex}
\newif\ifxetexorluatex
\ifxetex\xetexorluatextrue\fi
\ifluatex\xetexorluatextrue\fi

\ifxetexorluatex
    \usepackage{fontspec}
    \newfontfamily\georgiafont{TeX Gyre Pagella}   % sustituto libre de Georgia
    \newfontfamily\cambriafont{TeX Gyre Termes}    % sustituto libre de Cambria
    \newcommand{\georgiab}{\georgiafont\bfseries}
    \newcommand{\georgiamd}{\georgiafont}
    \newcommand{\cambriamd}{\cambriafont}
\else
    \newcommand{\georgiab}{\rmfamily\bfseries}
    \newcommand{\georgiamd}{\rmfamily\mdseries}
    \newcommand{\cambriamd}{\rmfamily\mdseries}
\fi

%========================================================================%
\makeatletter
\def\autor#1{\author{#1}}
\def\titulo#1{\title{#1}}
\def\fecha#1{\date{#1}}
\def\universidad#1{\gdef\@universidad{#1}}
\def\escuela#1{\gdef\@escuela{#1}}
\def\master#1{\gdef\@master{#1}}
\def\tipoTesis#1{\gdef\@tipoTesis{#1}}
\def\director#1{\gdef\@director{#1}}
\def\ciudad#1{\gdef\@ciudad{#1}}
\gdef\keywords#1{\gdef\keywordsv{#1}}
\gdef\keywordsEs#1{\gdef\keywordsESv{#1}}
\makeatother

\def\n{\bfseries}
\def\glmtr{\guillemotright}
\def\glmtl{\guillemotleft}

\newcolumntype{L}[1]{>{\hsize=#1\hsize\raggedright\arraybackslash}X}%
\newcolumntype{R}[1]{>{\hsize=#1\hsize\raggedleft\arraybackslash}X}%
\newcolumntype{C}[1]{>{\hsize=#1\hsize\centering\arraybackslash}X}%
\newcommand{\cc}[1]{\begin{minipage}[c][0.5cm][c]{0.5cm}\centering#1\end{minipage}}

\newenvironment{codigo}{
    \begin{minipage}[t][][t]{0.7\textwidth}
    \ttfamily \small \linespread{0}
}{
    \end{minipage}
}
\newcommand{\cod}{\ttfamily \small \linespread{0}}

\newenvironment{abstract}[1][Resumen]{%
    \begin{center}%
        {\fontsize{14}{14}\selectfont\bfseries #1\vspace{-0.5em}}%
    \end{center}%
    \quotation
}{\endquotation\vspace{1cm}}

%========================================================================%
\hypersetup{
    pdfpagemode={UseOutlines},
    bookmarksopen=true,
    bookmarksopenlevel=0,
    hypertexnames=false,
    colorlinks=true,
    pdfborder={0 0 0},
    citecolor=blue,
    linkcolor=blue,
    urlcolor=blue,
    pdfstartview={FitV},
    breaklinks=true
}

% Encabezado
\fancypagestyle{mainmatter}{%
    \fancyhead{}%
    \lhead{{\color{azul}Actividad grupal}}%
    \chead{}%
    \rhead{\masterv}%
    \lfoot{{\color{azul}Computación neuromórfica}}%
    \cfoot{}%
    \rfoot{\thepage\ de \getpagerefnumber{LastPage}}}

\fancypagestyle{frontmatter}{%
    \fancyhead{}%
    \lhead{{\color{azul}Actividad grupal}}%
    \chead{}%
    \rhead{\masterv}%
    \lfoot{{\color{azul}Computación neuromórfica}}%
    \cfoot{}%
    \rfoot{\thepage}}

%========================================================================%
\renewcommand*{\maketitle}{%
    \begin{titlepage}%
        {\flushleft%
            {\includegraphics[width=8.18cm]{logo-UNIR.pdf}}\par
        }%
        \vspace{1cm}
        \begin{table}[H]
        \hyphenpenalty=10000
        \setlength\arrayrulewidth{2pt}
            \begin{tabularx}{14.5cm}{ L{0.074} !{\color{azul}\vline}  L{0.931} }
                & \\
                & \fontsize{14}{14}\selectfont \georgiab \universidadv\\
                & \\
                & \fontsize{18}{18}\selectfont \georgiab \escuelav \\
                & \\
                & \fontsize{14}{14}\selectfont  \georgiab\masterv \\
                & \\ & \\ & \\
                & \fontsize{42}{42}\selectfont \cambriamd \color{azul}   \titulov \\
            \end{tabularx}
        \end{table}
        \vspace{0.5cm}
        {\flushleft \fontsize{11}{11}\selectfont %
            \georgiab \tipoTesisv\par
                \vspace{0.5cm}
                Presentado por: \georgiamd \autorv \par
                \vspace{0.5cm}
                \georgiab Director: \georgiamd \directorv \par
                \vspace{3cm}
                {\color{azul}%
                \georgiab Ciudad: \georgiamd \ciudadv \par
                \georgiab Fecha: \georgiamd \fechav \par}
        }
    \end{titlepage}}

%========================================================================%
% COMPACTADO DEL DOCUMENTO
% La clase 'book' abre cada capítulo en página nueva y reserva 50 pt sobre
% el título. En una memoria corta eso desperdicia media docena de páginas.
% Estas tres líneas eliminan el salto y reducen el espacio del encabezado,
% sin tocar tipografía ni márgenes: el texto se lee igual, solo que seguido.
% Para volver al comportamiento original de la plantilla, coméntalas.
\makeatletter
%\patchcmd{\chapter}{\if@openright\cleardoublepage\else\clearpage\fi}{}{}{}
% Espacios del encabezado de capítulo. Los tres valores originales de la
% clase 'book' son 50 pt encima, 20 pt entre "Capítulo N" y el título, y
% 40 pt debajo: 110 pt por capítulo, más de seis líneas de texto.
\patchcmd{\@makechapterhead}{50\p@}{6\p@}{}{}
\patchcmd{\@makechapterhead}{20\p@}{8\p@}{}{}
\patchcmd{\@makechapterhead}{40\p@}{12\p@}{}{}
% Evita líneas sueltas al principio o al final de página. No es infinito
% a propósito: con 10000 LaTeX prefiere dejar huecos grandes antes que
% romper, y eso gastaría más páginas de las que ahorra.
\widowpenalty=600
\clubpenalty=600
\displaywidowpenalty=600
\makeatother

%========================================================================%
% NOTAS AL EQUIPO. Se ven en rojo en el PDF para no perderlas de vista.
% Al resolver cada una, bórrala. Antes de entregar no debe quedar ninguna.
% Flotantes más pegados al texto: con [H] ya no se desplazan, así que el
% espacio por defecto de la clase 'book' resulta excesivo.
\setlength{\intextsep}{10pt plus 2pt minus 2pt}
\setlength{\textfloatsep}{10pt plus 2pt minus 2pt}

\newcommand{\equipo}[1]{{\color{red}\textbf{[EQUIPO]} \textit{#1}}}

%========================================================================%
\titulo{¿Qué retos plantean los nuevos modelos de computación?}
\autor{Benjamín David Salín Monroy, Héctor Lobato Silva, Nora Arroyo Carranza, Josué Isaac Gómez Méndez, Fernando Molina Mejía}
\director{Dr. Miguel Alejandro Medina Armendariz}
\fecha{12-10-2026}
\ciudad{CDMX}
\universidad{Universidad Internacional de la Rioja (UNIR)}
\escuela{Escuela Superior de Ingeniería y Tecnología}
\master{Máster en Computación Cuántica}
\tipoTesis{Actividad grupal · Computación neuromórfica}

\keywords{Neuromorphic computing, spiking neural networks, technology ethics}
\keywordsEs{Computación neuromórfica, redes de impulsos, ética de la tecnología}


%========================================================================%
\makeatletter
    \xdef\autorv{\@author}
    \xdef\titulov{\@title}
    \xdef\masterv{\@master}
    \xdef\fechav{\@date}
    \xdef\universidadv{\@universidad}
    \xdef\escuelav{\@escuela}
    \xdef\tipoTesisv{\@tipoTesis}
    \xdef\directorv{\@director}
    \xdef\ciudadv{\@ciudad}
\makeatother
\hypersetup{
    pdftitle={\titulov},
    pdfauthor={\autorv},
    pdfsubject={\masterv},
    pdfkeywords={\keywordsESv},
}

%========================================================================%
%                        INICIO DEL DOCUMENTO
%========================================================================%
\begin{document}

    %--- Páginas iniciales: portada, índice, resumen, listas ---%
    \frontmatter
    \pagestyle{frontmatter}
        \maketitle

        \begin{abstract}[Resumen]
            Este trabajo analiza los retos que plantea la computación
            neuromórfica, un paradigma que sitúa memoria y cómputo en el mismo
            punto físico y comunica mediante impulsos discretos dirigidos por
            eventos, frente a la arquitectura de von Neumann. El análisis se
            desarrolla en tres planos. En el tecnológico se identifican la
            ausencia de un modelo de programación común, la dificultad de
            entrenar redes de impulsos sin una función de activación derivable
            y la variabilidad entre dispositivos fabricados. En el ético se
            examinan la opacidad de los sistemas que aprenden tras su
            despliegue, la vigilancia ubicua que su bajo consumo hace viable y
            los riesgos asociados a las interfaces neuronales. En el social se
            valoran la concentración geográfica de la capacidad de fabricación,
            el doble uso civil y militar, y la paradoja de sostenibilidad que
            supone depender de materiales críticos. El equipo recopila las
            cifras medidas en cuatro dominios de aplicación, las somete a una
            lectura crítica que identifica sesgo de selección, fuentes no
            revisadas por pares y ausencia de un banco de pruebas común, y
            concluye que las medidas de mitigación deben incorporarse durante
            el diseño, no después del despliegue.
        \end{abstract}

%========================================================================%
%  CONTRIBUCIONES DEL EQUIPO
%  Es la evidencia del Criterio 3 (trabajo en equipo, 40 % de la nota).
%  Va sin numerar y antes de las referencias, como la declaración de
%  contribuciones de autoría de cualquier artículo científico.
%========================================================================%
{\large\bfseries\color{azul} Contribuciones del equipo}
\addcontentsline{toc}{chapter}{Contribuciones del equipo}
\vspace{0.6em}

El trabajo se desarrolló de forma colaborativa sobre un repositorio Git
compartido por los cinco integrantes. El historial de confirmaciones registra
qué modificó cada uno y en qué momento, y constituye la evidencia verificable de
la participación individual. La tabla \ref{tab:contribuciones} resume el reparto
de responsabilidades; la búsqueda de fuentes, la discusión de las conclusiones y
la revisión final se realizaron conjuntamente.

% Esta tabla NO es un flotante: como entorno 'table' no cabía al final de
% su página y dejaba media hoja en blanco. Compuesta como tabular normal,
% el texto fluye y se recupera esa página.
\begin{center}
\captionof{table}{Reparto de responsabilidades por apartado.}
\label{tab:contribuciones}
\small
\begin{spacing}{1}
\begin{tabularx}{\textwidth}{@{}L{0.62} L{0.95} L{0.43}@{}}
\toprule
\textbf{Integrante} & \textbf{Aportación} & \textbf{Secciones} \\
\midrule
Benjamín David Salín Monroy
  & Plataformas de hardware neuromórfico y dominios de aplicación con
    resultados medidos
  & \ref{sec:aplicaciones} \\
Fernando Molina Mejía
  & Plataformas de hardware neuromórfico y dominios de aplicación con
    resultados medidos
  & \ref{sec:aplicaciones} \\
Héctor Lobato Silva
  & Implicaciones éticas: opacidad y responsabilidad, vigilancia,
    neuroderechos, y paradoja de la sostenibilidad
  & \ref{sec:etica}, \ref{sec:impacto} \\
Josué Isaac Gómez Méndez
  & Impacto social: doble uso civil y militar, concentración tecnológica
    y brecha digital
  & \ref{sec:impacto} \\
Nora Arroyo Carranza
  & Estructura y maquetación, fundamentos teóricos, lectura crítica de
    las cifras y estrategias de mitigación
  & \ref{cap:introduccion}, \ref{sec:critica}, \ref{sec:desafios},
    \ref{sec:estrategias} \\
\midrule
Todo el equipo
  & Búsqueda y verificación de fuentes, conclusiones, postura del equipo
    y revisión final
  & \ref{cap:conclusiones} \\
\bottomrule
\end{tabularx}
\end{spacing}
\end{center}

% sloppypar permite espaciado más holgado en este párrafo: la URL no se
% puede partir y sin esto desborda el margen tres puntos.
\begin{sloppypar}
\noindent Repositorio del trabajo:
\url{https://github.com/HectorLobatoSilva/maestria_cuantica}, carpeta
\textit{Matematicas de la informacion / Actividad 2}.
\end{sloppypar}



        \clearpage
        % Los índices van con interlineado sencillo para que quepan en una
        % sola hoja; el cuerpo del documento conserva el 1,5 de la plantilla.
        \begin{spacing}{1}
            \small
            \tableofcontents
            % \listoffigures abriría capítulo y forzaría página nueva; se
            % compone a mano para que el índice de figuras quede en la misma
            % hoja que el índice general.
            \vspace{1.6em}
            {\large\bfseries\color{azul} Índice de figuras}\par\vspace{0.7em}
            \makeatletter\@starttoc{lof}\makeatother
        \end{spacing}



    %========================================================================%
    %  CUERPO DEL DOCUMENTO — ESTRUCTURA OBLIGATORIA UNIR
    %  (introducción, desarrollo, resultados, discusión, conclusiones,
    %   bibliografía). Escriba el contenido de cada sección directamente
    %  aquí abajo; ya no hace falta ningún \include ni carpeta "contenido/".
    %========================================================================%
    \mainmatter
    \pagestyle{mainmatter}

\chapter{Introducción}
\label{cap:introduccion}


Durante más de siete décadas la práctica totalidad de los sistemas de cómputo se
ha construido sobre la arquitectura de von Neumann, que separa la unidad de
procesamiento de la memoria. Esa separación obliga a transferir datos de forma
continua entre ambas, y ese tránsito, no la velocidad de cálculo, se ha
convertido en el principal límite de rendimiento y consumo; las cargas de
inteligencia artificial lo agravan, porque mueven volúmenes de parámetros muy
superiores. La computación neuromórfica propone una alternativa estructural:
hardware inspirado en la organización del sistema nervioso, donde procesamiento
y almacenamiento coinciden físicamente y la comunicación se produce mediante
impulsos discretos, solo cuando hay actividad \parencite{mead1990}. No se trata
de ejecutar más rápido el mismo modelo de cómputo, sino de reorganizarlo.

\begin{figure}[H]
\centering
\begin{tikzpicture}[font=\footnotesize,
    panel/.style={draw=gray!55,rounded corners=3pt,
                  minimum width=5.6cm,minimum height=1.8cm},
    caja/.style={draw=gray,fill=white,rounded corners=2pt,
                 minimum width=1.9cm,minimum height=0.8cm,align=center},
    celda/.style={draw=azul,fill=azul!12,
                  minimum width=1.85cm,minimum height=0.8cm,align=center}]
  % --- izquierda: arquitectura convencional
  \node[panel] at (0,0) {};
  \node[anchor=south] at (0,1.02) {\textbf{Arquitectura convencional}};
  \node[caja] (cpu) at (-1.45,0.25) {Procesador};
  \node[caja] (mem) at ( 1.45,0.25) {Memoria};
  \draw[<->,line width=1.6pt,red!65!black] (cpu) -- (mem);
  \node[color=red!65!black] at (0,-0.55) {el dato viaja en \emph{cada} operación};
  % --- derecha: arquitectura neuromórfica
  \node[panel] at (7.4,0) {};
  \node[anchor=south] at (7.4,1.02) {\textbf{Arquitectura neuromórfica}};
  \node[celda] at (6.48,0.25) {Cómputo};
  \node[celda,fill=azul!25] at (8.32,0.25) {Memoria};
  \draw[azul,line width=1.1pt,rounded corners=2pt]
        (5.52,-0.26) rectangle (9.28,0.76);
  \node[color=azul] at (7.4,-0.55) {el dato ya está donde se opera};
\end{tikzpicture}
\caption{Separación entre procesamiento y memoria frente a su colocalización.}
\label{fig:arquitectura}
\end{figure}

Su relevancia responde a dos presiones convergentes: el coste energético de
entrenar y ejecutar modelos de inteligencia artificial, que crece a un ritmo
difícil de sostener, y la necesidad de procesamiento autónomo en el propio
dispositivo, sin conexión permanente con centros de datos remotos, en ámbitos
como la robótica, los vehículos autónomos o las prótesis médicas
\parencite{roy2019}.

\section{Objetivo}

El objetivo es analizar los retos de la computación neuromórfica en tres planos:
tecnológico, ético y social.

Conviene cuantificarlo. Según \textcite{horowitz2014}, a 45 nanómetros una
operación aritmética o un acceso a la caché interna consumen unos 10 picojulios,
frente a los 1 o 2 nanojulios de un acceso a la DRAM externa: mover el dato es
unas cien veces más caro que operarlo. Esa proporción es el verdadero límite del
cómputo actual.

% El capítulo 2 arranca en hoja nueva. Sin esto, TeX reserva al final de la
% hoja anterior el hueco que necesitaría el encabezado del capítulo, y deja
% suelta la última línea del capítulo 1. Forzando el salto aquí, el
% capítulo 1 cierra entero en su hoja y no se gasta ninguna página extra.
\clearpage
\chapter{Desarrollo}
\label{cap:desarrollo}

\section{¿En qué consiste el paradigma?}

La computación neuromórfica diseña circuitos electrónicos que reproducen
principios de organización del tejido nervioso. Sus tres rasgos distintivos
frente a la arquitectura convencional son los siguientes.

\begin{itemize}
  \item \textbf{Colocalización de memoria y cómputo.} Las sinapsis artificiales
  almacenan los pesos en el mismo punto donde se realiza la operación, lo que
  elimina el tránsito constante de datos entre procesador y memoria.
  \item \textbf{Comunicación por impulsos.} La información viaja en forma de
  impulsos discretos, o \emph{spikes}, en lugar de valores continuos.
  \item \textbf{Operación dirigida por eventos.} No existe un reloj global que
  active todo el sistema en cada ciclo: una unidad consume energía únicamente
  cuando recibe o emite un impulso. De ahí procede la mayor parte del ahorro
  energético del paradigma \parencite{indiveri2015}.
\end{itemize}

\section{Fundamentos teóricos}

El punto de partida conceptual se encuentra en los sistemas VLSI analógicos
propuestos por \textcite{mead1990}, que reproducían en silicio la dinámica de las
neuronas biológicas aprovechando el comportamiento físico de los transistores en
régimen de subumbral.

El modelo de cómputo asociado son las redes neuronales de impulsos. En ellas, cada
neurona acumula las entradas que recibe a lo largo del tiempo y emite un impulso
cuando su potencial interno supera un umbral, tras lo cual el potencial se
reinicia. El modelo más utilizado es el de integración y disparo con fugas, cuya
dinámica puede expresarse como

\begin{equation}
\tau_m \frac{dV(t)}{dt} = -\bigl(V(t) - V_{\mathrm{reposo}}\bigr) + R\, I(t),
\label{eq:lif}
\end{equation}

donde $V(t)$ es el potencial de membrana, $\tau_m$ la constante de tiempo, $R$ la
resistencia de membrana e $I(t)$ la corriente de entrada. Cuando $V(t)$ alcanza
el umbral $V_{\mathrm{umbral}}$ se genera un impulso y el potencial vuelve a su
valor de reposo.

\begin{figure}[H]
\centering
\begin{tikzpicture}[font=\footnotesize]
  \draw[->,gray] (0,0) -- (9.4,0) node[right,color=gray]{$t$};
  \draw[->,gray] (0,-0.1) -- (0,2.3) node[above,color=gray]{$V$};
  \draw[dashed,red!60!black] (0,1.7) -- (9.2,1.7);
  \node[color=red!60!black,anchor=south west,font=\footnotesize] at (0.1,1.74) {umbral};
  \draw[azul,line width=1.1pt]
    (0,0.25) -- (1.1,0.85) -- (1.15,0.7) -- (2.2,1.4) -- (2.25,1.25)
    -- (3.1,1.71) -- (3.12,0.25) -- (4.3,0.8) -- (4.35,0.65)
    -- (5.4,1.2) -- (5.45,1.05) -- (6.6,1.71) -- (6.62,0.25)
    -- (8.0,0.75) -- (8.05,0.63) -- (9.0,1.1);
  \draw[line width=1.3pt,azul] (3.12,1.95) -- (3.12,2.4);
  \draw[line width=1.3pt,azul] (6.62,1.95) -- (6.62,2.4);
  \node[color=azul,anchor=west,font=\footnotesize] at (6.85,2.25) {impulsos emitidos};
\end{tikzpicture}
\caption{Potencial de membrana a lo largo del tiempo. Al alcanzar el umbral se
emite un impulso y el potencial se reinicia.}
\label{fig:neurona}
\end{figure}

La diferencia esencial con las redes neuronales artificiales convencionales es que
aquí el tiempo forma parte del cómputo: la información no se codifica solo en la
intensidad de una señal, sino en el instante y la frecuencia con que se producen
los impulsos.

\section{Aplicaciones prácticas y dominios de despliegue}
\label{sec:aplicaciones}

El paradigma ha salido del laboratorio y existe en silicio. Entre las
plataformas más consolidadas figuran \textbf{TrueNorth} (IBM), un circuito
digital de un millón de neuronas orientado a inferencia de muy bajo consumo
\parencite{merolla2014}; \textbf{Loihi} (Intel), que incorpora aprendizaje en el
propio chip mediante reglas de plasticidad sináptica \parencite{davies2018};
\textbf{SpiNNaker} (Universidad de Mánchester), una arquitectura masivamente
paralela para simular redes a gran escala en tiempo real \parencite{furber2014}.

Los ámbitos más explorados son el reconocimiento de patrones sensoriales con
restricciones de energía, la robótica autónoma, el procesamiento en el borde de
la red y las prótesis neuronales \parencite{schuman2022}. Una enumeración
genérica no permite valorar la madurez real del paradigma, así que a
continuación se examinan cuatro dominios con el resultado publicado que los
respalda.

\subsection{Navegación autónoma y visión basada en eventos}
Las cámaras de eventos generan datos asíncronos solo cuando cambia la
luminosidad de un píxel. Combinadas con redes de impulsos, el sistema procesa el
flujo óptico de forma nativa, sin barridos de fotogramas. Este enfoque reduce
hasta un 20\% el error de estimación frente a redes convencionales de tamaño
equivalente, con una décima parte de la energía y una fracción de los parámetros
\parencite{chowdhury2025}. Es el dominio con mejores resultados comparativos, y
el dato es doble: no solo ahorra energía, también mejora la precisión.

\subsection{Sensores acústicos siempre activos}
El reconocimiento de palabras clave en dispositivos móviles y de Internet de las
Cosas exige operación permanente sin agotar la batería. El procesador
\textit{Xylo Audio 2} realiza esa detección con 219 microvatios en reposo y
0,0066 milijulios por inferencia, manteniendo un acierto del 95,3\%
\parencite{bos2024}, lo que supera en un orden de magnitud a los
microcontroladores convencionales.

\subsection{Escalabilidad en centros de datos}
A escala de infraestructura, el objetivo es frenar el crecimiento del consumo
eléctrico del aprendizaje profundo. \textit{Hala Point}, que integra más de un
millar de chips Loihi 2 con 1.150 millones de neuronas y 128.000 millones de
sinapsis, procesa cargas cognitivas complejas con un consumo máximo de 2.600
vatios y una eficiencia superior a 15 TOPS/W \parencite{intel2024}.

\subsection{Neuroprótesis e interfaces cerebro-computadora}
Es el dominio de menor madurez clínica y mayor potencial. La capacidad de las
redes de impulsos para imitar la transmisión nerviosa permite concebir
interfaces que decodifiquen patrones de activación y los traduzcan en órdenes
para extremidades biónicas. Conviene afirmarlo con claridad: \textcite{donati2024}
concluyen que estas tecnologías \enquote{todavía tienen un largo camino por
recorrer} y que persisten retos de biocompatibilidad, estabilidad de la señal y
ausencia de marcos normativos. Lo que la revisión identifica son líneas de
investigación abiertas, no sistemas clínicos desplegados.

\subsection{Inteligencia artificial en el borde de la red}
La integración con el procesamiento en el borde permite que ciertos modelos
residan en el propio hardware, sin comunicación continua con la nube, lo que
reduce la dependencia de la infraestructura remota \parencite{schuman2022} y la
latencia, porque la información se procesa dirigida por eventos en lugar de
esperar a un reloj global \parencite{indiveri2015}. La magnitud del ahorro, eso
sí, conviene tomarla de las mediciones del apartado siguiente y no de las
estimaciones de orden de magnitud que circulan sin respaldo experimental.

\section{Lectura crítica de las cifras}
\label{sec:critica}

La tabla \ref{tab:cifras} reúne los cuatro dominios anteriores para los que
existe una cifra medida y publicada. Su elaboración es aportación propia del
equipo a partir de las cuatro fuentes citadas.

\begin{table}[H]
\centering
\caption{Resultados medidos por dominio de aplicación.}
\label{tab:cifras}
\begin{tabularx}{\textwidth}{L{0.8} L{0.7} L{1.5}}
\toprule
\textbf{Dominio} & \textbf{Plataforma} & \textbf{Dato medido} \\
\midrule
Visión autónoma & SNN frente a ANN & $-20$\% error; unas 10 veces menos energía \\
Audio siempre activo & Xylo Audio 2 & 219 $\mu$W en reposo; 0,0066 mJ por inferencia \\
Centro de datos & Hala Point & $1{,}15 \cdot 10^{9}$ neuronas en 2.600 W \\
Neuroprótesis & --- & sin cifras clínicas publicadas \\
\bottomrule
\end{tabularx}
\end{table}

Estas cifras exigen tres advertencias. Primera, hay \textbf{sesgo de selección}:
ninguna fuente compara el paradigma en una tarea donde salga perdiendo, de modo
que los datos son correctos pero no representan al conjunto de aplicaciones
posibles. Segunda, importa el \textbf{origen de la fuente}: la afirmación de que
los sistemas basados en Loihi operan con cien veces menos energía procede de una
nota de prensa del fabricante \parencite{intel2024}, no de una publicación
revisada por pares, y debe citarse como tal; las otras tres cifras sí proceden de
revistas con revisión. Tercera, \textbf{no existe un banco de pruebas común}: que
Xylo supere a Loihi y a SpiNNaker en energía por inferencia no significa que sea
superior, sino que resuelven problemas distintos sin una medida estándar que
permita compararlos. Esta carencia enlaza con el primer desafío tecnológico del
capítulo siguiente.

Existe además literatura independiente que matiza directamente estas cifras.
\textcite{yan2024} señalan que la mayoría de los trabajos sobre redes de
impulsos estiman el consumo contando únicamente las sumas realizadas, e ignoran
el coste de los accesos a memoria y del movimiento de datos. Al incorporarlos,
la ventaja energética deja de ser incondicional y pasa a depender de la ventana
temporal empleada, del nivel de dispersión de la actividad neuronal y de la
arquitectura concreta: para una red VGG16 con una ventana de seis pasos, los
autores calculan que la dispersión debe superar aproximadamente el 93\% para que
la red de impulsos ahorre energía en la mayoría de arquitecturas. Conviene
advertir que se trata de una preimpresión y no de un artículo revisado por
pares, pero su objeción es metodológica y afecta por igual a las cuatro cifras
de la tabla \ref{tab:cifras}: ninguna de las fuentes citadas descuenta el coste
del movimiento de datos. La ventaja del paradigma es real, pero está acotada a
las condiciones en las que se mide.

\chapter{Resultados}
\label{cap:resultados}

Este capítulo recoge el análisis del equipo en sus dos vertientes: los desafíos
técnicos que hoy frenan la adopción del paradigma y las implicaciones éticas que
se derivan de sus características propias.

\section{Desafíos tecnológicos}
\label{sec:desafios}

\subsection{Ausencia de un modelo de programación maduro}

La computación convencional dispone de un modelo de referencia estable y décadas
de herramientas construidas sobre él. La neuromórfica carece de equivalente: cada
plataforma impone su entorno, sus primitivas y sus limitaciones, de modo que el
código escrito para una no es trasladable a otra. Esa fragmentación eleva el
coste de entrada y dificulta comparar resultados entre grupos
\parencite{schuman2022}.

\subsection{Entrenamiento de las redes de impulsos}

La retropropagación, base del aprendizaje profundo, exige una función de
activación derivable, y el impulso no lo es: es un evento discreto. El
entrenamiento recurre entonces a aproximaciones —sustituir la derivada por una
función sustituta, convertir redes ya entrenadas o emplear reglas de plasticidad
local—, y ninguna alcanza todavía, de forma general, la precisión de los modelos
convencionales en tareas equivalentes \parencite{roy2019}.

\subsection{Fabricación y variabilidad de los dispositivos}

Las implementaciones analógicas y las basadas en memristores presentan dispersión
entre unidades del mismo proceso de fabricación, además de sensibilidad a la
temperatura y al envejecimiento. Dos chips nominalmente idénticos pueden dar
resultados distintos, lo que compromete la reproducibilidad y obliga a calibrar o
a tolerar la variación \parencite{indiveri2015}.

\section{Implicaciones éticas}
\label{sec:etica}

De los dilemas que plantea el paradigma, el equipo ha seleccionado los tres que
se derivan de forma más directa de sus características técnicas propias, y no de
la inteligencia artificial en general.

\subsection{Opacidad, dilución de la responsabilidad y pérdida de control humano}
Un sistema neuromórfico aprende y modifica su comportamiento de forma continua, distribuida y local tras su despliegue gracias a su capacidad de aprendizaje \textit{in situ}, que en plataformas como Loihi se implementa mediante reglas de plasticidad sináptica ejecutadas en el propio chip \parencite{davies2018}. Al carecer de un conjunto de parámetros fijos e inspeccionables, resulta sumamente difícil reconstruir el razonamiento detrás de una decisión concreta y atribuir responsabilidades cuando el sistema falla. Esta \enquote{opacidad extrema} agrava los retos éticos tradicionales de la inteligencia artificial, haciendo que la mitigación de sesgos algorítmicos y el mantenimiento del control humano sean técnicamente más complejos que en los modelos convencionales de redes neuronales.

La consecuencia jurídica es inmediata. Si uno de estos sistemas, integrado en
infraestructuras críticas, robótica médica o vehículos autónomos, toma una
decisión perjudicial tras haber modificado su propio comportamiento de forma
impredecible, resulta legal y éticamente complejo determinar quién responde del
fallo: ni el fabricante entregó ese comportamiento, ni el operador lo programó.
La rendición de cuentas se diluye precisamente porque la adaptación autónoma sin
supervisión constante es una característica deseada del diseño, no un defecto.

\subsection{Vigilancia ubicua y la paradoja del bajo consumo}
El bajo consumo energético y la velocidad de ejecución son las principales virtudes técnicas del paradigma neuromórfico \parencite{schuman2022}, pero paradójicamente constituyen el origen de su riesgo social más inmediato. Su extrema eficiencia permite instalar sensores capaces de analizar imagen y sonido de forma permanente, alimentados por baterías de larga duración o mediante la captación de energía ambiental. Al superarse la viabilidad económica de una vigilancia continua en el espacio público, esta deja de ser una barrera técnica para convertirse puramente en una decisión política, lo que exige nuevos marcos éticos para prevenir el control poblacional.

\subsection{La frontera con lo humano, la biomimética y los neuroderechos}
La computación neuromórfica difumina la línea entre la máquina y el ser humano al implementar copias analógicas o digitales de las redes neuronales biológicas en circuitos electrónicos, utilizando elementos como el memristor para simular sinapsis \parencite{indiveri2015}. Al describirse mediante vocabulario neurológico y aplicarse a interfaces cerebro-computadora (BCI) o prótesis conectadas al sistema nervioso, el paradigma reabre debates fundamentales sobre la autonomía personal, la agencia y la privacidad mental. Ante esta convergencia, marcos globales y expertos coinciden en la necesidad de establecer prioridades éticas específicas para las neurotecnologías, protegiendo así la identidad individual, el libre albedrío y los datos neuronales de los usuarios \parencite{yuste2017}.

\chapter{Discusión}
\label{cap:discusion}

\section{Estrategias para abordar los riesgos éticos}
\label{sec:estrategias}

Las tres medidas que proponemos se corresponden, una a una, con las tres
implicaciones expuestas en el capítulo anterior: la auditabilidad responde a la
opacidad, la regulación del despliegue a la vigilancia ubicua y la gobernanza
específica a los riesgos de las interfaces neuronales.

\subsection{Auditabilidad por diseño}

Frente a la opacidad, la medida no consiste en renunciar al aprendizaje continuo,
sino en exigir que todo sistema desplegado en un ámbito sensible registre de forma
verificable las entradas recibidas, los cambios aplicados a sus parámetros y las
decisiones emitidas. Ese registro debe estar disponible para auditores externos y
conservarse durante un plazo definido. La condición ha de incorporarse en la fase
de diseño, porque añadirla después resulta técnicamente inviable.

\subsection{Regulación del despliegue en espacios públicos}

Frente a la vigilancia ubicua, la respuesta corresponde al ámbito normativo:
delimitar en qué lugares y con qué finalidades puede instalarse un sensor con
capacidad de análisis permanente, establecer la obligación de señalizarlo y
prohibir el reconocimiento biométrico indiscriminado. Marcos como el Reglamento
europeo de inteligencia artificial ofrecen un punto de partida que puede adaptarse
a otros contextos regulatorios.

\subsection{Gobernanza específica de la neurotecnología}

Frente a los riesgos asociados a las interfaces neuronales, proponemos exigir
revisión por comités éticos independientes, consentimiento informado que
contemple de manera explícita el uso posterior de los datos neuronales, y un
tratamiento jurídico diferenciado de esos datos, que no pueden equipararse a
cualquier otro dato personal \parencite{yuste2017}.

\section{Impacto en la sociedad y en la humanidad}
\label{sec:impacto}

El impacto social de este paradigma admite lecturas opuestas, y conviene
recorrerlas en orden.

\textbf{Los beneficios.} Al procesar de una forma más cercana al funcionamiento
del cerebro y evitar el movimiento constante de información entre memoria y
procesador \parencite{roy2019}, el paradigma permite ejecutar modelos
directamente en el dispositivo. Hoy buena parte de las aplicaciones de
inteligencia artificial dependen de enviar información a servidores remotos; la
computación neuromórfica reduce esa dependencia \parencite{schuman2022}, con el
consiguiente efecto ambiental positivo. En el ámbito médico abre la puerta a
prótesis inteligentes, interfaces neuronales y dispositivos capaces de analizar
señales biológicas en tiempo real durante periodos prolongados, lo que mejoraría
la autonomía de personas con alguna discapacidad. Ahora bien, cuanto más directo
sea el contacto con el sistema nervioso, más necesario será proteger esa
información y la identidad de quien la genera \parencite{yuste2017}.

\textbf{La vigilancia.} Un dispositivo de bajo consumo capaz de analizar señales
de forma permanente es también una herramienta de monitoreo continuo. El
problema se agrava al incorporar reconocimiento facial, identificación
biométrica o algoritmos capaces de inferir comportamientos y relaciones entre
personas: puede existir observación permanente del entorno sin que nadie sea
consciente de ello.

\textbf{El doble uso.} El contexto de los conflictos armados lo hace evidente. El
Comité Internacional de la Cruz Roja advierte que el uso de inteligencia
artificial en el ámbito militar puede aumentar la velocidad y la escala de los
conflictos y reducir el control humano sobre las decisiones relativas al uso de
la fuerza \parencite{icrc2026ai}. El mismo procesador que permite a una prótesis
analizar señales durante horas puede mantener un dron o una red de sensores
operando durante días para identificar movimientos, objetos o personas. La
autonomía, por tanto, no se evalúa de forma aislada: depende de quién utiliza el
sistema y con qué finalidad. En un dispositivo médico es una ventaja; en uno
militar o de vigilancia, permitir que una máquina clasifique personas sin
supervisión suficiente tiene consecuencias de otro orden.

\textbf{La concentración.} Diseñar procesadores avanzados exige capital,
conocimiento e infraestructura especializada. La OCDE documenta que la capacidad
mundial de fabricación de obleas está muy concentrada geográficamente y que la
escasa sustituibilidad entre plantas convierte esa concentración en un cuello de
botella capaz de amplificar cualquier interrupción del suministro
\parencite{oecd2025chips}. La UNCTAD añade que el acceso a la infraestructura y a
la experiencia en inteligencia artificial sigue concentrado en unas pocas
economías, y lo cuantifica: apenas cien empresas, radicadas sobre todo en Estados
Unidos y China, concentran el 40\% del gasto empresarial mundial en investigación
y desarrollo \parencite{unctad2025}. La brecha tecnológica ya no consiste en
quién accede a una computadora, sino en quién puede diseñar los procesadores y
controlar la infraestructura que los ejecuta.

\textbf{La sostenibilidad.} Aunque el paradigma busca reducir el coste
energético de los centros de datos \parencite{roy2019}, su fabricación depende de
dispositivos memristivos y materiales no convencionales. El problema no
desaparece: se desplaza del consumo eléctrico hacia la extracción de recursos
críticos y la gestión de los residuos. En 2022 se generaron 62 millones de
toneladas de residuos electrónicos, de las que solo el 22,3\% se trató de forma
documentada, y el reciclaje cubre como máximo el 1\% de la demanda mundial de
tierras raras esenciales \parencite{ewaste2024}. Un paradigma que depende de esos
materiales no puede presentarse como sostenible atendiendo solo a su consumo en
operación.

En conjunto, el impacto dependerá menos de la tecnología que de quién la
controle y para qué se use. Las mismas características que favorecen
aplicaciones médicas y sistemas más eficientes sirven para construir vigilancia
permanente, sistemas militares más autónomos y nuevas formas de recopilación de
información. El reto no es solo diseñar procesadores más eficientes, sino
establecer los mecanismos de transparencia, regulación y responsabilidad que
permitan aprovecharlos sin ampliar al mismo tiempo la capacidad de vigilancia,
control o desigualdad.

\section{Percepción de la humanidad y evolución tecnológica}

Imitar la arquitectura del cerebro no equivale a reproducir lo que el cerebro
hace. Conviene mantener la distinción, porque el vocabulario del paradigma
(neuronas, sinapsis, aprendizaje) induce a atribuir a estos sistemas capacidades
de las que carecen.

La metáfora, además, se ha invertido. Durante décadas el cerebro se explicaba
como un ordenador; ahora se construyen ordenadores que se explican como
cerebros. El riesgo de esa inversión no es técnico sino interpretativo: cuanto
mejor funcione un sistema descrito con vocabulario neurológico, mayor será la
tentación de concluir que entre procesar información y comprenderla solo media
una diferencia de grado. Nada de lo expuesto en esta memoria lo respalda.

Sobre la evolución tecnológica, el paradigma apunta a un cambio de dirección más
que de velocidad. Durante medio siglo el progreso se midió en transistores por
unidad de superficie; agotado ese escalado, la mejora depende de rediseñar la
arquitectura en lugar de miniaturizarla. La computación neuromórfica es uno de
esos rediseños, igual que la cuántica, y ambas comparten un rasgo revelador:
ninguna aspira a sustituir al modelo convencional, sino a convivir con él
resolviendo aquello en lo que este resulta ineficiente. El escenario más
probable no es la sucesión de paradigmas, sino la fragmentación del cómputo en
arquitecturas especializadas. Esa fragmentación tiene un coste: multiplica los
modelos de programación, dificulta comparar sistemas y concentra en todavía
menos actores la capacidad de diseñarlos, con los efectos descritos en la
sección \ref{sec:impacto}.

\chapter{Conclusiones}
\label{cap:conclusiones}

El análisis permite al equipo formular cuatro conclusiones.

\textbf{Primera.} El paradigma no acelera el modelo convencional: lo reorganiza.
Al unir memoria y cómputo y operar por eventos elimina el cuello de botella de la
arquitectura de von Neumann, a costa de renunciar a su generalidad.

\textbf{Segunda.} Sus obstáculos ya no son de fabricación sino de programación y
entrenamiento: existe hardware funcional, pero falta un modelo de programación
común y un método de aprendizaje que iguale de forma general a los convencionales.

\textbf{Tercera.} Las implicaciones éticas no son un añadido externo: proceden de
su mayor virtud técnica. Un consumo tan reducido es, a la vez, lo que permite
desplegar sensores permanentes en cualquier lugar.

\textbf{Cuarta.} Las medidas de mitigación deben incorporarse durante el diseño.
La auditabilidad no se añade a un sistema que ya aprende, y regular el despliegue
pierde eficacia una vez instalada la infraestructura.

\section*{Postura del equipo}

A la pregunta de si la computación neuromórfica sustituirá al modelo
convencional o convivirá con él, el equipo sostiene la segunda respuesta: será
un acelerador especializado, no un sustituto. Tres elementos de los capítulos
anteriores la respaldan.

\textbf{La naturaleza de su ventaja.} Las cifras de la tabla \ref{tab:cifras}
proceden de tareas sensoriales y dispersas: flujo óptico, palabra clave, señales
neuronales. Ninguna fuente lo compara en cómputo de propósito general, y
\textcite{yan2024} precisan el motivo: la ventaja energética depende de que la
actividad sea muy dispersa. Un paradigma cuya eficiencia está condicionada a una
propiedad del problema no puede reemplazar a uno que no impone esa condición.

\textbf{El estado de la programación.} Mientras no exista un modelo común ni un
método de entrenamiento equivalente \parencite{schuman2022, roy2019}, no puede
absorber la base instalada de software, mejore lo que mejore el hardware.

\textbf{El precedente.} La unidad de procesamiento gráfico no sustituyó al
procesador central: se especializó junto a él y redefinió qué cargas corresponden
a cada uno. Todo apunta a que este paradigma sigue esa trayectoria, ocupando el
procesamiento sensorial permanente y de bajo consumo en el borde de la red.

Esto no rebaja su importancia, la refuerza: si su lugar son los sensores
permanentes, las implicaciones éticas del capítulo \ref{cap:resultados} no son un
escenario hipotético sino su consecuencia directa.

    %--- Referencias (formato APA, al final del documento) ---%
    %  \clearpage para que arranquen en hoja propia y no compartan página
    %  con las conclusiones.
    \clearpage
    \chapter*{Referencias}
    \addcontentsline{toc}{chapter}{Referencias}
    % Interlineado sencillo y sin separación extra entre entradas, para que
    % las 18 referencias quepan en una sola hoja.
    % OJO: biblatex ignora un \scriptsize suelto delante del entorno; el
    % cuerpo de letra de la bibliografía solo se cambia con \bibfont.
    \renewcommand*{\bibfont}{\footnotesize}
    \begin{spacing}{1}
        \setlength{\bibitemsep}{0pt plus 0pt minus 0pt}
        % Permite que las URLs largas (CICR, UNESCO, UNCTAD) se partan entre
        % líneas. Sin esto, una URL que no cabe desborda el margen y genera
        % el aviso "Overfull \hbox" en esta línea.
        \setcounter{biburlnumpenalty}{9000}
        \setcounter{biburlucpenalty}{9000}
        \setcounter{biburllcpenalty}{9000}
        \printbibliography[heading=none]
    \end{spacing}

\end{document}
